\section{Related work}
\label{sec:related}
\noindent\textbf{Visual Question Answering (VQA)} is the task of responding to natural language queries by interpreting visual content~\cite{antol2015vqa,Naik_2023_ICCV,Li_2024_CVPR}. Recent approaches leverage vision-language models (VLMs) built on transformer architectures~\citep{fields2023vision, shi2024non, xiao2022video, lei2021less, cheng2023vindlu} alongside pre-trained language models~\citep{yang2022zero, yu2023self}. For example, \citet{cho2021unifying} introduced a generative transformer model that integrates visual and textual modalities for VQA. Many approaches enhance generalization by leveraging compositionality—a cornerstone of cognitive reasoning~\citep{keysers2019measuring, lake2017building}. For instance, \citet{johnson2017clevr} explored the composition of visual attributes by creating a dataset designed for compositional reasoning, while \citet{whitehead2021separating} used contrastive learning to enhance compositionality, disentangling reasoning skills from visual concepts. Despite these advances, implicit decomposition strategies can hinder generalization, and crafting effective contrastive samples remains challenging.
 
\noindent \textbf{Continual Learning (CL)} seeks to develop frameworks that incrementally assimilate new data while mitigating catastrophic forgetting. This is a fundamental challenge for many deep learning methods due to catastrophic forgetting (CF)~\citep{mcclelland1995there}. CL methods are broadly categorized into \emph{knowledge distillation}-based approaches, which constrain mappings between successive models to prevent forgetting~\citep{li2017learning, icarl, 8953661, 9880426, douillard2020podnetpooledoutputsdistillation,Kang2022afc}; \emph{optimisation}-based, which adjust gradient updates to minimize interference with prior tasks~\cite{lopez2017gradient,chaudhry2018efficient, saha2021gradient, Yang_2023_ICCV} and \emph{representation}-based strategies that learn robust, adaptable features~\cite{gao2023unified,foret2020sharpness,ermis2022memory,douillard2022dytox}. Recently, \emph{prompt}-based methods have emerged~\citep{wang2022learning, wang2022dualprompt, wang2022s, smith2023coda}, employing visual prompts with pre-trained transformers in CL scenarios~\citep{liu2023pre}. However, since most techniques are designed for Class-Incremental Learning (CIL) in unimodal settings, they fall short when applied to multimodal data and overlook the critical issue of compositional generalization in VQA~\citep{zhang2023vqacl, nikandrou2022task}.

\noindent \textbf{Continual VQA. }Recent studies have explored multimodal continual learning in VQA~\citep{del2020ratt, greco2019psycholinguistics, nikandrou2022task, srinivasan2022climb}. However, prior works like \citet{greco2019psycholinguistics} primarily analyze forgetting dynamics without proposing dedicated solutions, often overlooking the role of pre-trained models~\citep{mehta2021empirical}. The VQACL benchmark~\citep{zhang2023vqacl} evaluated conventional continual learning approaches such as ~\citep{chaudhry2019tinyepisodicmemoriescontinual, Chaudhry2019er, Wan_2022_CVPR, zhang2023vqacl}. Other studies have investigated specific facets of continual VQA, including question diversity~\citep{greco2019psycholinguistics}, compositionality~\citep{zhang2023vqacl} and domain adaptation~\citep{zhang-etal-2022-continual}. In contrast, our work introduces a novel question-only replay mechanism with attention distillation for continual VQA that is both memory-efficient, and privacy-conscious, eliminating the need for stored prototypes as employed in the VQACL method~\citep{zhang2023vqacl}\footnote{VQACL represents both the setting and the approach for continual learning in Visual Question Answering, as described in ~\citep{zhang2023vqacl}.}. 